% \vspace{-0.4cm}
\section{Related Work}
% \vspace{-0.3cm}
Table~\ref{tab:overview} (Appendix~\ref{app:relwork-table}) surveys prior work across three related areas; we explain the key distinctions below.

\subsection{Agentic Data and Tool-Calling Benchmarks}
The dominant paradigm for evaluating tool-calling agents uses either environment-based execution feedback~\citep{zhou2024webarena,drouin2024workarena,trivedi2024appworld,yao2024taubench} or deterministic scoring against annotated trajectories~\citep{deng2023mind2web,xu2025agenttrek,patil2025bfcl}. $\tau$-bench~\citep{yao2024taubench} stresses agents against simulated users and proposes \texttt{pass\textasciicircum k} reliability rather than per-turn correctness; like this line generally, it scores agents but does not interrogate the judge producing those scores. BFCL~\citep{patil2025bfcl} is the closest structural ancestor to our data format, isolating single-turn function-calling with typed JSON schemas and AST-based ground-truth comparison; we adopt the same schema but extend it to multi-step DAG-structured workflows at three controlled difficulty tiers. The most structurally similar work is FuncBenchGen~\citep{maekawa2025funcbenchgen}, which frames multi-step calling as DAG traversal with controllable complexity. The key distinction is purpose: FuncBenchGen trains and evaluates \emph{generator} agents, whereas AgentJudgeBench measures \emph{judge} reliability on those generators' outputs via a paired with-GT/without-GT protocol and per-metric decomposition that FuncBenchGen does not provide. TaskBench~\citep{shen2024taskbench} shares the tool-dependency graph framing but focuses on decomposition quality rather than judge alignment. Synthetic data pipelines~\citep{wang2023selfinstruct,xu2025magpie,cui2024ultrafeedback} inform our generation approach but target model training rather than evaluation benchmarking.

\subsection{LLM-as-Judge}
\citet{zheng2023judging} established the LLM-as-judge paradigm on MT-Bench, showing GPT-4's strong agreement with human preferences on open-ended dialogue while identifying positional preference, verbosity sensitivity, and self-enhancement as systematic biases. Subsequent work either extends the evaluation surface -- JudgeBench~\citep{tan2025judgebench} and Arena-Hard-Auto~\citep{li2024arenahard} move to hard, verifiable response pairs -- or builds dedicated judge models~\citep{wang2024pandalm,li2024autoj,zhu2023judgelm,kim2024prometheus2}. JudgeLM~\citep{zhu2023judgelm} fine-tunes 7B--33B judges on GPT-4-distilled verdicts and identifies position, knowledge, and format biases that mirror those reported for prompt-based judges; Prometheus~2~\citep{kim2024prometheus2} adds rubric-conditioned direct assessment with open weights. Both target text-quality scoring rather than structural tool-calling correctness, but their bias taxonomies inform our prompt-format ablation (\S\ref{sec:abl:format}). The jury-of-judges paradigm~\citep{verga2024jury} shows that ensembling reduces individual judge biases. Crucially, this line targets text-centric tasks where judge quality reduces to preference alignment over fluent output. AgentJudgeBench occupies a different regime: correctness is structural, the verdict space is $\{0, 0.5, 1\}$ rather than a preference ranking, and failure modes are multi-dimensional and interdependent.

\subsection{LLM Judges for Tool-Calling}
LLM judges have been deployed for tool-calling evaluation in several recent benchmarks, but always as an implementation detail rather than the subject of study. ToolLLM~\citep{qin2024toolllm} uses ChatGPT to compute pass rates over API trajectories, reporting 87.1\% human agreement in aggregate; StableToolBench~\citep{guo2024stabletoolbench} and MCP-AgentBench~\citep{guo2025mcpagentbench} adopt similar approaches with updated judge models. GeoBenchX~\citep{krechetova2025geobenchx} assembles a three-judge panel achieving 88--96\% agreement, and Agent-as-a-Judge~\citep{zhuge2024agentjudge} and Auto-Eval Judge~\citep{bhonsle2025autoeval} propose more structured evaluation frameworks. In each case, judge reliability is reported as a single aggregate figure on a small sample, with no variation of task complexity, difficulty, or ground-truth availability. ToolSandbox~\citep{lu2024toolsandbox} notably questions LLM judge reliability directly but replaces judges with programmatic evaluation rather than characterising their failure modes. AgentJudgeBench takes the complementary stance: keep the judge paradigm and measure it systematically, decomposing reliability by metric, topology, difficulty, and condition across $321{,}648$ completed evaluations (Appendix~\ref{app:eval-stats}). A detailed feature comparison across all seven related systems appears in Appendix~\ref{app:judge-comparison}.
